\subsection{QUOTIENT LAWS}

DEFINITION 11. Let E be a set. A law of composition \(\top\) and an equivalence relation R on E are said to be compatible if the relations \(x \equiv x'\) (mod R) and \(y \equiv y'\) (mod R) (for \(x, x', y, y' \in \mathbf{E}\) ) imply \(x \top y \equiv x' \top y'\) (mod R); the law of composition on the quotient set E/R which maps the equivalence classes of \(x\) and \(y\) to the equivalence class of \(x \top y\) is called the quotient law of the law \(\top\) with respect to R. The set E/R with the quotient law is called the quotient magma of E with respect to R.

To say that an equivalence relation R on E is compatible with the internal law of composition \(f: E \times E \to E\) on E means that the mapping f is compatible (in the sense of Set Theory, II, §6, no. 5) with the product equivalence relation \(R \times R\) on \(E \times E\) and the equivalence relation R on E. (Set Theory, II, §6, no. 8). This also means that the graph of R is a submagma of \(E \times E\) .

If the law \(\top\) is associative (resp. commutative) so is the quotient law (more briefly we say that associativity, or commutativity, is preserved when passing to the quotient).

The canonical mapping from the magma E to the magma E/R is a homomorphism.

For a mapping g of E/R into a magma F to be a homomorphism it is necessary and sufficient that the composition of g with the canonical mapping of E onto E/R be a homomorphism.

The two following propositions are immediate from the definitions:

PROPOSITION 6. Let E and F be two magmas and f a homomorphism of E into F. Let R\{x,y\} denote the relation \(f(x)=f(y)\) between elements x,y of E. Then R is an equivalence relation on E compatible with the law on E and the mapping of E/R onto f(E) derived from f by passing to the quotient is an isomorphism of the quotient magma E/R onto the submagma f(E) of F.

PROPOSITION 7. Let E be a magma and R an equivalence relation on E compatible with the law on E. For an equivalence relation S on E/R to be compatible with the quotient law it is necessary and sufficient that S be of the form T/R where T is an equivalence relation on E implied by R and compatible with the law on E. The canonical mapping of E/T onto (E/R)/(T/R) (Set Theory, II, §6, no. 7) is then a magma isomorphism.

PROPOSITION 8. Let E be a magma, A a stable subset of E and R an equivalence relation on E compatible with the law on E. The saturation B of A with respect to R (Set

Theory, II, § 6, no. 5) is a stable subset. The equivalence relations \(\mathbf{R}_{\mathbf{A}}\) and \(\mathbf{R}_{\mathbf{B}}\) induced by \(\mathbf{R}\) on \(\mathbf{A}\) and \(\mathbf{B}\) respectively are compatible with the induced laws and the mapping derived from the canonical injection of \(\mathbf{A}\) into \(\mathbf{B}\) by passing to the quotients is a magma isomorphism of \(\mathbf{A} / \mathbf{R}_{\mathbf{A}}\) onto \(\mathbf{B} / \mathbf{R}_{\mathbf{B}}\) .

Let \(\top\) denote the law on E. If \(x\) and \(y\) are two elements of B there exist two elements \(x'\) and \(y'\) of A such that \(x \equiv x' \pmod{R}\) and \(y \equiv y' \pmod{R}\) ; then \(x \top y \equiv x' \top y' \pmod{R}\) and \(x' \top y' \in A\) , whence \(x \top y \in B\) . Thus B is a stable subset of E and the other assertions are obvious.

Let M be a magma and \(((u_{\alpha}, v_{\alpha}))_{\alpha \in I}\) a family of elements of \(M \times M\) . Consider all the equivalence relations S on M which are compatible with the law on M and such that \(u_{\alpha} \equiv v_{\alpha} \pmod{S}\) for all \(\alpha \in I\) . The intersection of the graphs of these relations is the graph of an equivalence relation R which is compatible with the law on M and such that \(u_{\alpha} \equiv v_{\alpha} \pmod{R}\) . Hence R is the finest (Set Theory, III, §1, nos. 3 and 7) equivalence relation with these two properties. It is called the equivalence relation compatible with the law on M generated by the \((u_{\alpha}, v_{\alpha})\) .

PROPOSITION 9. Preserving the above notation, let \(f\) be a homomorphism of \(\mathbf{M}\) into a magma such that \(f(u_{\alpha}) = f(v_{\alpha})\) for all \(\alpha \in \mathbf{I}\) . Then \(f\) is compatible with \(\mathbf{R}\) .

Let T be the equivalence relation associated with f.~Then \(u_{\alpha} \equiv v_{\alpha} \pmod{T}\) for all \(\alpha \in I\) and T is compatible with the law on M, hence T is coarser than R; this proves the proposition.

